\chapter{Concepts et technologies}
\label{ch:etatart}

\section{De l'ETL à la Modern Data Stack}

Les chaînes décisionnelles classiques reposent sur le modèle \textbf{ETL} (\emph{Extract, Transform,
Load}) : les données sont transformées par un outil dédié avant d'être chargées dans l'entrepôt. Avec
l'arrivée des entrepôts de données dans le cloud, dont la puissance de calcul est élastique, le modèle
\textbf{ELT} s'est imposé : les données sont d'abord chargées brutes dans l'entrepôt, puis transformées
\emph{dans} l'entrepôt, en SQL.

La \textbf{Modern Data Stack} désigne l'ensemble d'outils spécialisés et interopérables qui
accompagnent ce modèle :
\begin{itemize}
  \item un \textbf{entrepôt cloud} séparant stockage et calcul (Snowflake, BigQuery, Redshift\ldots) ;
  \item un outil d'\textbf{ingestion} déclaratif (Fivetran, Airbyte, DLT\ldots) ;
  \item un outil de \textbf{transformation} versionné et testé (dbt) ;
  \item un \textbf{orchestrateur} (Airflow, Dagster, Prefect) ;
  \item des outils de \textbf{restitution} (tableaux de bord, applications analytiques, IA).
\end{itemize}

Cette approche applique à la donnée les pratiques du génie logiciel : gestion de versions, tests
automatisés, revue de code, intégration continue et déploiement reproductible.

\section{L'architecture Médaillon}

L'architecture Médaillon, popularisée par Databricks, organise les données d'un \emph{lakehouse} ou d'un
entrepôt en trois couches de qualité croissante (figure~\ref{fig:medaillon}).

\begin{figure}[H]
\centering
\begin{tikzpicture}[node distance=1.3cm]
  \node[layer, fill=bronze] (b) {BRONZE\\[-1pt]\scriptsize\mdseries données brutes};
  \node[layer, fill=silver, right=of b] (s) {SILVER\\[-1pt]\scriptsize\mdseries données nettoyées};
  \node[layer, fill=gold, right=of s] (g) {GOLD\\[-1pt]\scriptsize\mdseries données métier};
  \draw[flow] (b) -- (s);
  \draw[flow] (s) -- (g);
  \node[below=0.25cm of b, text width=3.4cm, align=center, font=\scriptsize]
    {Copie fidèle de la source,\\ historique d'ingestion,\\ aucune règle métier};
  \node[below=0.25cm of s, text width=3.4cm, align=center, font=\scriptsize]
    {Typage, nettoyage,\\ dédoublonnage,\\ normalisation, conformité};
  \node[below=0.25cm of g, text width=3.4cm, align=center, font=\scriptsize]
    {Modèle en étoile,\\ indicateurs, agrégats,\\ prêt pour l'analyse et l'IA};
\end{tikzpicture}
\caption{Principe de l'architecture Médaillon}
\label{fig:medaillon}
\end{figure}

\begin{itemize}
  \item \textbf{Bronze} : les données sont chargées telles quelles depuis la source. Cette couche
        garantit la traçabilité et permet de rejouer les transformations sans réinterroger la source.
  \item \textbf{Silver} : les données sont nettoyées, typées, dédoublonnées et conformées. Elles restent
        proches de la structure source mais sont fiables.
  \item \textbf{Gold} : les données sont modélisées pour un usage métier --- ici un schéma en étoile ---
        et directement exploitables par les outils d'analyse et par l'IA.
\end{itemize}

\section{La modélisation dimensionnelle}

La couche Gold suit la \textbf{modélisation dimensionnelle} de Ralph Kimball, qui organise les données
autour :
\begin{itemize}
  \item de \textbf{tables de faits}, qui enregistrent des événements mesurables (une ligne de facture,
        un mouvement de stock) avec leurs mesures numériques (quantité, montant, bénéfice) ;
  \item de \textbf{tables de dimensions}, qui décrivent le contexte de ces événements (qui, quoi, où,
        quand) avec des attributs descriptifs (nom du client, catégorie, ville, mois\ldots).
\end{itemize}

Le \textbf{grain} d'une table de faits --- le niveau de détail d'une ligne --- est la décision de
conception centrale. Ce modèle simplifie fortement les requêtes analytiques : un indicateur s'obtient
par agrégation des faits et regroupement selon des attributs de dimensions, avec des jointures simples
et prévisibles. Cette propriété est précieuse pour la génération automatique de SQL par un LLM.

\subsection{Dimensions à évolution lente (SCD)}
Les attributs des dimensions changent dans le temps : un produit change de prix, un client de plafond de
crédit. Kimball propose plusieurs stratégies, dont :
\begin{itemize}
  \item \textbf{SCD de type 1} : l'ancienne valeur est écrasée (pas d'historique) ;
  \item \textbf{SCD de type 2} : chaque changement crée une nouvelle ligne, avec une période de
        validité (\texttt{valid\_from}, \texttt{valid\_to}) et un indicateur de version courante.
\end{itemize}

\section{Les outils de la plateforme}

\subsection{SQL Server}
SGBD relationnel de Microsoft, utilisé ici comme \textbf{système source} transactionnel hébergeant la
base WideWorldImporters (SQL Server 2025, édition Developer).

\subsection{Snowflake}
Snowflake est un entrepôt de données proposé en SaaS sur les principaux clouds. Il sépare le
\textbf{stockage} (facturé au volume) du \textbf{calcul}, assuré par des \emph{virtual warehouses}
dimensionnables, suspendus automatiquement lorsqu'ils sont inactifs. Il offre notamment :
\begin{itemize}
  \item un contrôle d'accès par rôles (RBAC) hiérarchiques ;
  \item le \textbf{clonage zéro copie} (\emph{zero-copy cloning}) d'une base entière, instantané et sans
        duplication physique des données ;
  \item la \emph{Time Travel} et des \emph{resource monitors} pour maîtriser la consommation ;
  \item \textbf{Snowflake Cortex}, une suite de services d'IA exécutés au sein de la plateforme, dont
        \textbf{Cortex Analyst} pour le Text-to-SQL ;
  \item \textbf{Streamlit in Snowflake}, l'hébergement d'applications Streamlit directement dans le
        compte, sans infrastructure à gérer.
\end{itemize}

\subsection{DLT (data load tool)}
DLT est une bibliothèque Python open source d'ingestion de données. Elle extrait les données d'une source
(ici une base SQL via SQLAlchemy), infère et fait évoluer le schéma cible, normalise les noms, et charge
les données dans la destination. Elle gère nativement le \textbf{chargement incrémental} (curseur sur
une colonne de date) et les modes d'écriture \texttt{append}, \texttt{replace} et \texttt{merge}
(fusion sur clé primaire), ainsi que la conservation de l'état entre deux exécutions.

\subsection{dbt (data build tool)}
dbt permet d'écrire les transformations sous forme de modèles \texttt{SELECT} en SQL enrichi de Jinja.
dbt résout les dépendances entre modèles (fonction \texttt{ref()}), les matérialise dans l'entrepôt
(vues, tables, tables incrémentales), exécute des \textbf{tests} de qualité (unicité, non-nullité,
intégrité référentielle, tests personnalisés), génère la \textbf{documentation} et le \textbf{lignage},
et propose des \textbf{snapshots} pour historiser des données en SCD de type 2.

\subsection{Dagster}
Dagster est un orchestrateur orienté \emph{assets} : au lieu de décrire des tâches, on déclare les
\textbf{actifs de données} produits (tables, modèles) et leurs dépendances. Dagster en déduit le graphe
d'exécution, offre une interface de visualisation du lignage, des planifications (\emph{schedules}), des
capteurs (\emph{sensors}), des politiques de reprise sur erreur, et des contrôles de qualité
(\emph{asset checks}). Ses intégrations \texttt{dagster-dlt} et \texttt{dagster-dbt} transforment
automatiquement chaque table DLT et chaque modèle dbt en asset, reliés dans un même lignage.

Le sujet laissait le choix entre Dagster et Airflow ; Dagster a été retenu car il est explicitement
recommandé par l'énoncé, que son approche par assets correspond naturellement à une architecture en
couches, et qu'il s'intègre nativement avec DLT et dbt.

\subsection{Streamlit}
Streamlit est un cadriciel Python permettant de construire rapidement des applications web de données
interactives, sans développement front-end : composants de saisie, tableaux, graphiques, interface de
conversation (\emph{chat}).

\section{Le Text-to-SQL et Snowflake Cortex Analyst}

Le \textbf{Text-to-SQL} consiste à traduire une question en langage naturel en une requête SQL
exécutable. Les grands modèles de langage ont fait progresser considérablement cette tâche, mais une
génération « naïve » --- envoyer la question et la liste des colonnes à un LLM --- reste fragile :
ambiguïté des termes métier (« ventes » : HT ou TTC ?), jointures erronées, colonnes inventées,
calculs incorrects.

\textbf{Cortex Analyst} est le service de Snowflake dédié à ce problème. Il s'appuie sur un
\textbf{modèle sémantique} (fichier YAML) qui décrit :
\begin{itemize}
  \item les tables logiques et leurs colonnes, classées en \emph{dimensions}, \emph{dimensions
        temporelles}, \emph{faits} et \emph{métriques} ;
  \item des descriptions et des \textbf{synonymes} métier (par exemple « CA » pour le chiffre d'affaires) ;
  \item les \textbf{relations} (jointures) entre tables ;
  \item des \textbf{requêtes vérifiées} (\emph{verified queries}) servant d'exemples de référence ;
  \item des instructions personnalisées (langue de réponse, conventions de calcul).
\end{itemize}

Cortex Analyst s'expose sous forme d'API REST : il reçoit la conversation et le modèle sémantique, et
renvoie une interprétation de la question, la requête SQL générée ou, en cas d'ambiguïté, des
suggestions de reformulation. Il ne fait qu'\emph{écrire} le SQL : son exécution reste sous le
contrôle de l'application, avec les droits de l'utilisateur.
